The following section presents studies of the METNet ML algorithm within the dark photon analysis.
Within ATLAS, \(E_T^{\text{miss}} \) is reconstructed using various working points.
These working points differ in the selection criteria of input jets, such as  $p_T$ and JVT tagging requirements.
Each analysis selects a WP based on performance in the relevant phase space and pileup rejection.
These WPs have a high dependence on event topology, and are not optimal for every event that might enter the signal region.
A new NN tool, METNet~\cite{hodkinson_2023}, has been developed to allow for WP selection on an event-by-event basis.

METNet classifies events based on Equation~\ref{eq:metnet}, and the outputs are N weights, \( \omega_N \), corresponding to the likelihood that each WP is optimal for the event.
The weights sum to unity, and are used to build a prediction for \(E_T^{\text{miss}} \) as shown in Equation~\ref{eq:metnet_etmiss}, summed in quadrature.

\begin{equation}
    Class = \text{argmin}_{WP} \left[ (p^{miss,WP}_{x} - p^{miss,true}_{x})^2 + (p^{miss,WP}_{y} - p^{miss,true}_{y})^2 \right]
    \label{eq:metnet}
\end{equation}
\begin{equation}
    p^{miss,METNet}_{x,y} = \sum_{WP} \omega_{WP} p^{miss,WP}_{x,y}
    \label{eq:metnet_etmiss}
\end{equation}

Performance studies of METNet have been conducted previously and show an increase in performance on a variety of event topologies~\cite{hodkinson_2023,Horstmann2025metnet}, such as $t\bar{t}$ and $Z\rightarrow\mu\mu$.
The CP task for the dark photon analysis is to validate the latest version of the METNet algorithm and to study the impact on our signal sensitivity.
The METNet algorithm has been integrated into the analysis framework, and studies are presented here using four MC samples shown in table \ref{tab:metnet_samples}.
A network file specific to Run 3 conditions was provided by the METNet developers.

\begin{table}[htbp]
    \centering
    \caption{MC samples used for METNet studies.}
    \begin{tabular}{l l}
        \toprule
        \hline
        Process & Sample Name  \\
        \hline
        $ggF\rightarrow H \rightarrow\gamma\gamma_D$ & \tiny\texttt{\detokenize{mc23_13p6TeV.602402.PhPy8EG_ggH_H125yyd.deriv.DAOD_PHYS.e8537_e8528_s4159_s4114_r15224_r15225_p6320}}  \\
        $\gamma$ + jets fragmentation & \tiny\texttt{\detokenize{mc23_13p6TeV.801649.Py8_gammajet_frag_DP8_17_FullSW.deriv.DAOD_PHYS.e8514_e8528_s4159_s4114_r15224_r15225_p6491}}  \\
         & \tiny\texttt{\detokenize{mc23_13p6TeV.801650.Py8_gammajet_frag_DP17_35_FullSW.deriv.DAOD_PHYS.e8514_e8528_s4159_s4114_r15224_r15225_p6491}} \\
         & \tiny\texttt{\detokenize{mc23_13p6TeV.801651.Py8_gammajet_frag_DP35_50_FullSW.deriv.DAOD_PHYS.e8514_e8528_s4159_s4114_r15224_r15225_p6491}} \\
         & \tiny\texttt{\detokenize{mc23_13p6TeV.801652.Py8_gammajet_frag_DP50_70_FullSW.deriv.DAOD_PHYS.e8514_e8528_s4159_s4114_r15224_r15225_p6491}} \\
         & \tiny\texttt{\detokenize{mc23_13p6TeV.801653.Py8_gammajet_frag_DP70_140_FullSW.deriv.DAOD_PHYS.e8514_e8528_s4159_s4114_r15224_r15225_p6491}} \\
         & \tiny\texttt{\detokenize{mc23_13p6TeV.801654.Py8_gammajet_frag_DP140_280_FullSW.deriv.DAOD_PHYS.e8514_e8528_s4159_s4114_r15224_r15225_p6491}} \\
         & \tiny\texttt{\detokenize{mc23_13p6TeV.801655.Py8_gammajet_frag_DP280_500_FullSW.deriv.DAOD_PHYS.e8514_e8528_s4159_s4114_r15224_r15225_p6491}} \\
         & \tiny\texttt{\detokenize{mc23_13p6TeV.801656.Py8_gammajet_frag_DP500_800_FullSW.deriv.DAOD_PHYS.e8514_e8528_s4159_s4114_r15224_r15225_p6491}} \\
         & \tiny\texttt{\detokenize{mc23_13p6TeV.801657.Py8_gammajet_frag_DP800_1000_FullSW.deriv.DAOD_PHYS.e8514_e8528_s4159_s4114_r15224_r15225_p6491}} \\
         & \tiny\texttt{\detokenize{mc23_13p6TeV.801658.Py8_gammajet_frag_DP1000_1500_FullSW.deriv.DAOD_PHYS.e8514_e8528_s4159_s4114_r15224_r15225_p6491}} \\
         & \tiny\texttt{\detokenize{mc23_13p6TeV.801659.Py8_gammajet_frag_DP1500_2000_FullSW.deriv.DAOD_PHYS.e8514_e8528_s4159_s4114_r15224_r15225_p6491}} \\
         & \tiny\texttt{\detokenize{mc23_13p6TeV.801660.Py8_gammajet_frag_DP2000_inf_FullSW.deriv.DAOD_PHYS.e8514_e8528_s4159_s4114_r15224_r15225_p6491}} \\
        $Z\rightarrow\mu\mu$ + jets & \tiny\texttt{\detokenize{mc23_13p6TeV.700789.Sh_2214_Zmumu_maxHTpTV2_BFilter.deriv.DAOD_PHYS.e8514_e8528_s4159_s4114_r15224_r15225_p6491}} \\
         & \tiny\texttt{\detokenize{mc23_13p6TeV.700790.Sh_2214_Zmumu_maxHTpTV2_CFilterBVeto.deriv.DAOD_PHYS.e8514_e8528_s4159_s4114_r15224_r15225_p6491}} \\
         & \tiny\texttt{\detokenize{mc23_13p6TeV.700791.Sh_2214_Zmumu_maxHTpTV2_CVetoBVeto.deriv.DAOD_PHYS.e8514_e8528_s4159_s4114_r15530_r15514_p6491}} \\
        $Z\rightarrow\nu\nu+\gamma$ & \tiny\texttt{\detokenize{mc23_13p6TeV.700776.Sh_2214_nunugamma.deriv.DAOD_PHYS.e8514_e8528_s4159_s4114_r15224_r15225_p6491}} \\
        \hline
        \bottomrule
    \end{tabular}
    \label{tab:metnet_samples}
\end{table}

Besides the dark photon signal sample, three background samples are used for METNet studies: $\gamma$ + jets fragmentation, $Z\rightarrow\mu\mu$ + jets, and $Z\rightarrow\nu\nu+\gamma$.
$Z\rightarrow\mu\mu$ + jets was chosen as a cross check for the performance studies conducted in \cite{hodkinson_2023,Horstmann2025metnet}, while $Z\rightarrow\nu\nu+\gamma$ provides a true \(E_T^{\text{miss}} \) environment.
The $\gamma$ + jets fragmentation background enters the SR due to large amount of fake \(E_T^{\text{miss}} \) from primary vertex misidentification, making it a useful test of METNet performance in a challenging \(E_T^{\text{miss}} \) environment.

The cross check of METNet performance using $Z\rightarrow\mu\mu$ + jets is shown in Figure~\ref{fig:metnet_zmumu} without any analysis selections, with the RMS of the \(E_T^{\text{miss}} \) in x and y used as the metric for resolution.

\begin{figure}[htbp]
    \centering
    \begin{minipage}[b]{0.45\textwidth}
        \includegraphics[width=\textwidth]{../images/dark_photon/images/metnet/MET_zmumu_jets_netgood.pdf}
        \caption*{(a) \(E_T^{\text{miss}} \)}
    \end{minipage}
    \hfill
    \begin{minipage}[b]{0.45\textwidth}
        \includegraphics[width=\textwidth]{../images/dark_photon/images/metnet/rms_hist_vs_npv_netgood_zmumu_jets_netgood.pdf}
        \caption*{(b) RMS}
    \end{minipage}
    \caption{\(E_T^{\text{miss}} \) distribution (a) for METNet (blue), the tight WP (red), and truth (green) for $Z\rightarrow\mu\mu$ + jets. The resolution (b) is shown for all events in solid blue rectangles (METNet) and red triangles (tight WP), and plus marks and asterisks for good and bad reconstructed vertices respectively.}\label{fig:metnet_zmumu}
\end{figure}

As expected, METNet more closely matches the truth \(E_T^{\text{miss}} \) distribution compared to the tight WP.
The resolution as a function of $N_{\text{PV}}$ outperforms the tight WP until reaching a high number of vertices.
Interestingly, the difference in resolution for events with a bad reconstructed vertex is much larger for METNet compared to the tight WP. 
This indicates that METNet is sensitive to the correct identification of the primary vertex.
Overall, the results agree almost exactly with the previous studies~\cite{hodkinson_2023,Horstmann2025metnet}, confirming the proper integration of METNet into the analysis framework.

The \(E_T^{\text{miss}} \) distributions for the four samples are shown in Figure~\ref{fig:metnet_etmiss} with the analysis preselections applied.

\begin{figure}[htbp]
    \centering
    \begin{minipage}[b]{0.45\textwidth}
        \includegraphics[width=\textwidth]{../images/dark_photon/images/metnet/MET_ggHyyd_selections_netgood.pdf}
        \caption*{(a) ggF $H\rightarrow\gamma\gamma_D$}
    \end{minipage}
    \hfill
    \begin{minipage}[b]{0.45\textwidth}
        \includegraphics[width=\textwidth]{../images/dark_photon/images/metnet/MET_gamma_jet_selections_netgood.pdf}
        \caption*{(b) $\gamma$ + jets fragmentation}
    \end{minipage}

    \vspace{0.4cm}
    
    \begin{minipage}[b]{0.45\textwidth}
        \includegraphics[width=\textwidth]{../images/dark_photon/images/metnet/MET_zmumu_jets_selections_netgood.pdf}
        \caption*{(c) $Z\rightarrow\mu\mu$ + jets}
    \end{minipage}\
    \hfill
    \begin{minipage}[b]{0.45\textwidth}
        \includegraphics[width=\textwidth]{../images/dark_photon/images/metnet/MET_znunu_gamma_selections_netgood.pdf}
        \caption*{(d) $Z\rightarrow\nu\nu+\gamma$}
    \end{minipage}

    \caption{\(E_T^{\text{miss}} \) distributions for METNet (blue), the tight WP (red), and truth (green) for: (a) ggF $H\rightarrow\gamma\gamma_D$, (b) $\gamma$ + jets fragmentation, (c) $Z\rightarrow\mu\mu$ + jets, and (d) $Z\rightarrow\nu\nu+\gamma$. Analysis preselections applied.}\label{fig:metnet_etmiss}
\end{figure}

The METNet distributions trend slightly lower than the tight WP for all samples, with overall agreement in shape between the two.
While this brings METNet marginally more in line with truth for signal and $\gamma$ + jets, it appears to shift slightly lower than truth for $Z\rightarrow\nu\nu+\gamma$.
Note that the analysis trigger removes almost all $Z\rightarrow\mu\mu$ + jets events, rendering the sample less useful for comparison.
Resolutions with preselections applied are shown in Figure~\ref{fig:metnet_resolution}.

\begin{figure}[htbp]
    \centering
    \begin{minipage}[b]{0.45\textwidth}
        \includegraphics[width=\textwidth]{../images/dark_photon/images/metnet/rms_hist_vs_npv_netgood_ggHyyd_selections_netgood.pdf}
        \caption*{(a) ggF $H\rightarrow\gamma\gamma_D$}
    \end{minipage}
    \hfill
    \begin{minipage}[b]{0.45\textwidth}
        \includegraphics[width=\textwidth]{../images/dark_photon/images/metnet/rms_hist_vs_npv_netgood_gamma_jet_selections_netgood.pdf}
        \caption*{(b) $\gamma$ + jets fragmentation}
    \end{minipage}

    \vspace{0.4cm}
    
    \begin{minipage}[b]{0.45\textwidth}
        \includegraphics[width=\textwidth]{../images/dark_photon/images/metnet/rms_hist_vs_npv_netgood_zmumu_jets_selections_netgood.pdf}
        \caption*{(c) $Z\rightarrow\mu\mu$ + jets}
    \end{minipage}\
    \hfill
    \begin{minipage}[b]{0.45\textwidth}
        \includegraphics[width=\textwidth]{../images/dark_photon/images/metnet/rms_hist_vs_npv_netgood_znunu_gamma_selections_netgood.pdf}
        \caption*{(d) $Z\rightarrow\nu\nu+\gamma$}
    \end{minipage}

    \caption{Resolution as a function of $N_{\text{PV}}$ for METNet (blue) and the tight WP (red) for: (a) ggF $H\rightarrow\gamma\gamma_D$, (b) $\gamma$ + jets fragmentation, (c) $Z\rightarrow\mu\mu$ + jets, and (d) $Z\rightarrow\nu\nu+\gamma$. Analysis preselections applied.}\label{fig:metnet_resolution}
\end{figure}

The resolutions are generally comparable between METNet and the tight WP across all samples except for $Z\rightarrow\nu\nu+\gamma$.
As seen before, METNet degrades in performance at high $N_{\text{PV}}$ for this sample, and is outperformed by the tight WP at an $N_{\text{PV}}$ of ~25.
This could be due to the tight WP being specifically optimized for a true \(E_T^{\text{miss}} \) environment, while METNet may weight the other WPs too heavily for this topology.
Most importantly, there is no siginificant increase in resolution for signal events using METNet.
Acceptances for signal and the $\gamma$ + jets fragmentation background were calculated in comparison to the 100 GeV cut on \(E_T^{\text{miss}} \) used in the analysis.
It was found that a 100 GeV cut on METNet achieves the same 1.43\% acceptance for signal (1.71x$10^{-4}$\% for $\gamma$ + jets fragmentation), giving no gain in efficiency.

In summary, these studies have validated METNet performance on $Z\rightarrow\mu\mu$ + jets events, in agreement with previous studies.
When applying analysis preselections, METNet shows comparable performance to the tight WP in \(E_T^{\text{miss}} \) resolution across four samples, with some degradation at high pileup for $Z\rightarrow\nu\nu+\gamma$ events.
There is no significant improvement in resolution or signal acceptance using METNet compared to the tight WP, and the current recommendation of the tight WP is maintained.
However, the dark photon analysis presents a challenging environment for \(E_T^{\text{miss}} \) reconstruction due to the high rate of fake \(E_T^{\text{miss}} \) from primary vertex misidentification.
It is recommended that METNet be trained more extensively on events with low track activity and fake \(E_T^{\text{miss}} \), and revisited in the next iteration of the analysis.